% ****************************************************************************************************
% classicthesis-config.tex 
% Thèse Rémi Paccou — CIRED — 2026
% ****************************************************************************************************

\PassOptionsToPackage{utf8}{inputenc}
	\usepackage{inputenc}

% ****************************************************************************************************
% 1. Configure classicthesis
% ****************************************************************************************************
\PassOptionsToPackage{eulerchapternumbers,listings,%
					 pdfspacing,%
					 subfig,beramono,eulermath,parts,dottedtoc}{classicthesis}                                        

% ****************************************************************************************************
% 2. Personal data
% ****************************************************************************************************
\newcommand{\myTitle}{Information et calcul dans l'économie de l'énergie\xspace}
\newcommand{\mySubtitle}{Histoire, dynamiques économiques et l'hypothèse d'un couplage structurel\xspace}
\newcommand{\myDegree}{Docteur en Sciences Économiques\xspace}
\newcommand{\myName}{Rémi Paccou\xspace}
\newcommand{\myProf}{Directeur à compléter\xspace}
\newcommand{\myOtherProf}{Co-directeur à compléter\xspace}
\newcommand{\mySupervisor}{Superviseur à compléter\xspace}
\newcommand{\myFaculty}{AgroParisTech\xspace}
\newcommand{\myDepartment}{Sciences Économiques\xspace}
\newcommand{\myUni}{Université Grenoble Alpes\xspace}
\newcommand{\myLocation}{Grenoble\xspace}
\newcommand{\myTime}{2026\xspace}
\newcommand{\myVersion}{version de travail\xspace}

\newcounter{dummy}
\newlength{\abcd}
\providecommand{\mLyX}{L\kern-.1667em\lower.25em\hbox{Y}\kern-.125emX\@}
\newcommand{\ie}{i.\,e.}
\newcommand{\Ie}{I.\,e.}
\newcommand{\eg}{e.\,g.}
\newcommand{\Eg}{E.\,g.} 

% ****************************************************************************************************
% 3. Packages
% ****************************************************************************************************
\PassOptionsToPackage{french}{babel}
	\usepackage{babel}                  

\usepackage{csquotes}
\PassOptionsToPackage{%
    backend=biber,%
	language=auto,%
    style=authoryear-comp,%
    bibstyle=authoryear,dashed=false,%
    sorting=nyt,%
    maxbibnames=10,%
    natbib=true%
}{biblatex}
    \usepackage{biblatex}

\PassOptionsToPackage{fleqn}{amsmath}
    \usepackage{amsmath}

\PassOptionsToPackage{T1}{fontenc}
    \usepackage{fontenc}     
\usepackage{textcomp}
\usepackage{scrhack}
\usepackage{xspace}
\usepackage{mparhack}
\usepackage{fixltx2e}
\PassOptionsToPackage{printonlyused,smaller}{acronym} 
    \usepackage{acronym}
    \renewcommand*{\aclabelfont}[1]{\acsfont{#1}}

% ****************************************************************************************************
% 4. Floats
% ****************************************************************************************************
\usepackage{tabularx}
    \setlength{\extrarowheight}{3pt}
\newcommand{\tableheadline}[1]{\multicolumn{1}{c}{\spacedlowsmallcaps{#1}}}
\newcommand{\myfloatalign}{\centering}
\usepackage{caption}
\captionsetup{font=small}
\usepackage{subfig}  

% ****************************************************************************************************
% 5. Listings
% ****************************************************************************************************
\usepackage{listings} 
\lstset{language=[LaTeX]Tex,%
    morekeywords={PassOptionsToPackage,selectlanguage},
    keywordstyle=\color{RoyalBlue},
    basicstyle=\small\ttfamily,
    commentstyle=\color{Green}\ttfamily,
    stringstyle=\rmfamily,
    numbers=none,
    numberstyle=\scriptsize,
    stepnumber=5,
    numbersep=8pt,
    showstringspaces=false,
    breaklines=true,
    belowcaptionskip=.75\baselineskip
} 

% ****************************************************************************************************
% 6. Hyperreferences
% ****************************************************************************************************
\PassOptionsToPackage{pdftex,hyperfootnotes=false,pdfpagelabels}{hyperref}
    \usepackage{hyperref}
\pdfcompresslevel=9
\pdfadjustspacing=1 
\PassOptionsToPackage{pdftex}{graphicx}
    \usepackage{graphicx} 

\hypersetup{%
    colorlinks=true, linktocpage=true, pdfstartpage=3, pdfstartview=FitV,%
    breaklinks=true, pdfpagemode=UseNone, pageanchor=true, pdfpagemode=UseOutlines,%
    plainpages=false, bookmarksnumbered, bookmarksopen=true, bookmarksopenlevel=1,%
    hypertexnames=true, pdfhighlight=/O,%
    urlcolor=webbrown, linkcolor=RoyalBlue, citecolor=webgreen,%
    pdftitle={\myTitle},%
    pdfauthor={\textcopyright\ \myName, \myUni, \myFaculty},%
    pdfsubject={\mySubtitle},%
    pdfkeywords={ICT, information, calcul, énergie, économie, climat, datacenter, GPT},%
    pdfcreator={pdfLaTeX},%
    pdfproducer={LaTeX with hyperref and classicthesis}%
}   

\makeatletter
\@ifpackageloaded{babel}%
    {%
       \addto\extrasfrench{%
			\renewcommand*{\figureautorefname}{Figure}%
			\renewcommand*{\tableautorefname}{Tableau}%
			\renewcommand*{\partautorefname}{Partie}%
			\renewcommand*{\chapterautorefname}{Chapitre}%
			\renewcommand*{\sectionautorefname}{Section}%
			\renewcommand*{\subsectionautorefname}{Section}%
			\renewcommand*{\subsubsectionautorefname}{Section}%     
                }%
            \providecommand{\subfigureautorefname}{\figureautorefname}%             
    }{\relax}
\makeatother

% ****************************************************************************************************
% 7. Last calls
% ****************************************************************************************************
\listfiles
\usepackage{classicthesis} 

% --- Fix francisation des titres (classicthesis force "Contents" en anglais) ---
\addto\captionsfrench{\renewcommand{\contentsname}{Sommaire}}
\addto\captionsfrench{\renewcommand{\listfigurename}{Table des figures}}
\addto\captionsfrench{\renewcommand{\listtablename}{Liste des tableaux}}

% ****************************************************************************************************
% 8. Marges style Pottier (zone de texte étroite, marges généreuses)
% ****************************************************************************************************
\usepackage{booktabs}
\usepackage{enumitem}

\areaset[current]{420pt}{770pt}
\setlength{\marginparwidth}{7em}
\setlength{\marginparsep}{1em}

% ****************************************************************************************************
% 9. Commandes de travail (RETIRER EN VERSION FINALE)
% ****************************************************************************************************
\usepackage{xcolor}
\usepackage{svg}
\usepackage{tikz}
\definecolor{emergence}{rgb}{0.0, 0.45, 0.0}

\newcommand{\emergence}[1]{%
  \marginpar{\scriptsize\color{emergence}\textbf{É}}%
  {\color{emergence}\textbf{[Émergence]} #1}}

\newcommand{\questionch}[2]{%
  \marginpar{\scriptsize\color{RoyalBlue}\textbf{$\rightarrow$#1}}%
  {\color{RoyalBlue}\textit{[$\rightarrow$Ch.#1] #2}}}

\newcommand{\wip}[1]{#1}
% ****************************************************************************************************

% ****************************************************************************************************
% 10. Code couleur des écoles économiques
% ----------------------------------------------------------------------------------------------------
% Convention : marquer uniquement la sentence fondamentale (l'argument central)
% de chaque école, pas tout le texte qui en relève. Une phrase, parfois deux.
% Distinct des marqueurs de travail : teal (questions ouvertes), red (lacunes),
% blue (décisions architecturales).
% ****************************************************************************************************
\definecolor{schoolNeo}{RGB}{30,60,140}        % NEO    bleu marine
\definecolor{schoolInfo}{RGB}{110,60,160}      % INFO   violet
\definecolor{schoolAus}{RGB}{220,120,30}       % AUS    orange
\definecolor{schoolEvo}{RGB}{50,130,70}        % EVO    vert forêt
\definecolor{schoolInst}{RGB}{125,80,40}       % INST   sépia
\definecolor{schoolReg}{RGB}{140,30,55}        % REG    bordeaux
\definecolor{schoolSts}{RGB}{160,130,50}       % STS    ocre
\definecolor{schoolMarx}{RGB}{180,40,30}       % MARX   rouge cinabre

\newcommand{\NEO}[1]{\textcolor{schoolNeo}{\textbf{#1}}}     % Néoclassique standard
\newcommand{\INFO}[1]{\textcolor{schoolInfo}{\textbf{#1}}}    % Économie de l'information
\newcommand{\AUS}[1]{\textcolor{schoolAus}{\textbf{#1}}}      % Autrichienne / Hayekienne
\newcommand{\EVO}[1]{\textcolor{schoolEvo}{\textbf{#1}}}      % Évolutionniste-schumpétérienne
\newcommand{\INST}[1]{\textcolor{schoolInst}{\textbf{#1}}}    % Néo-institutionnaliste
\newcommand{\REG}[1]{\textcolor{schoolReg}{\textbf{#1}}}      % Post-keynésienne / Régulationniste
\newcommand{\STS}[1]{\textcolor{schoolSts}{\textbf{#1}}}      % Sociotechnique / LTS
\newcommand{\MARX}[1]{\textcolor{schoolMarx}{\textbf{#1}}}    % Marxiste (classique, contemporain, écologique, numérique)

% Marqueurs de marge légers, signalent qu'une école est mobilisée dans la section/paragraphe
% sans surcolorier le texte. Complémentaires aux macros inline ci-dessus.
\newcommand{\markNEO}{\marginpar{\scriptsize\color{schoolNeo}\textbf{$\alpha$\,NEO}}}
\newcommand{\markINFO}{\marginpar{\scriptsize\color{schoolInfo}\textbf{$\beta$\,INFO}}}
\newcommand{\markAUS}{\marginpar{\scriptsize\color{schoolAus}\textbf{$\gamma$\,AUS}}}
\newcommand{\markEVO}{\marginpar{\scriptsize\color{schoolEvo}\textbf{$\delta$\,EVO}}}
\newcommand{\markINST}{\marginpar{\scriptsize\color{schoolInst}\textbf{$\varepsilon$\,INST}}}
\newcommand{\markREG}{\marginpar{\scriptsize\color{schoolReg}\textbf{$\zeta$\,REG}}}
\newcommand{\markSTS}{\marginpar{\scriptsize\color{schoolSts}\textbf{$\eta$\,STS}}}
\newcommand{\markMARX}{\marginpar{\scriptsize\color{schoolMarx}\textbf{$\theta$\,MARX}}}

% ****************************************************************************************************
% 11. Marqueurs de la grille faits forts / contradictions productives (idée Rémi)
% ----------------------------------------------------------------------------------------------------
% Convention : marquer la qualification épistémique d'un mécanisme dans la marge.
% \factfort       : convergence de plusieurs écoles épistémiquement indépendantes
% \contradiction  : les écoles désaccordent, le test va trancher (chapitre 4)
% \proposition    : mécanisme proposé sans triangulation, statut conjectural
% ****************************************************************************************************
\definecolor{factfortcol}{RGB}{20,90,40}
\definecolor{contradcol}{RGB}{180,90,30}
\definecolor{propocol}{RGB}{100,100,100}

\newcommand{\factfort}[1]{%
  \marginpar{\scriptsize\color{factfortcol}\textbf{$\checkmark\checkmark$}}%
  {\color{factfortcol}\textbf{[Fait fort]} #1}}

\newcommand{\contradiction}[1]{%
  \marginpar{\scriptsize\color{contradcol}\textbf{$\rightleftarrows$}}%
  {\color{contradcol}\textbf{[Contradiction productive]} #1}}

\newcommand{\proposition}[1]{%
  \marginpar{\scriptsize\color{propocol}\textbf{?}}%
  {\color{propocol}\textbf{[Proposition]} #1}}
% ****************************************************************************************************
